% IEEE journal source; author metadata remains to be supplied by the user.
% Full proofs are inlined. The external analytical plot is optional.
\documentclass[journal]{IEEEtran}
\usepackage[T1]{fontenc}
\usepackage{amsmath,amssymb,amsthm,booktabs,graphicx,hyperref}
\newtheorem{theorem}{Theorem}
\newtheorem{proposition}{Proposition}
\newtheorem{corollary}{Corollary}
\title{Diagnostic Information Limits under Healthy Drift and Closed-Loop Memory}
\author{Author information to be supplied}
\begin{document}
\maketitle
% Body fragment: uses the host document's theorem environments and packages.
\begin{abstract}
We study deterministic diagnostic calendars for a known stable linear plant
with complete retained states and a scalar healthy input that drifts throughout
holds and missing intervals. When the fixed Gaussian process-noise image
contains the fault direction, the constant-force score loss across a missing
span determines two regimes: a sharp deficit of order five halves when that loss
is positive, and matching quadratic-order bounds when the score is preserved.
Both results compare with an ideal healthy-subtracted complete record and use
the complete input clock, calendar-uniform lower bounds, and observable
attaining constructions. Positive-definite noise forces the first regime;
a separate compact common-noise controller theorem permits joint leading-order
design. Controller examples and an independent nonlinear mechanical certificate
explain the information restrictions and finite-risk obligations of these laws.
\end{abstract}

\section{Introduction}\label{sf:introduction}
Known feedback transforms the fault, healthy input and process noise through
the same plant. With complete state records, inverting that transformation
preserves the diagnostic experiment. A missing state can prevent this inversion,
so the diagnostic effect of a controller depends on which states are retained,
rather than on output noise correlation alone.

We consider two task signs, a growing known fault, and a separate unrestricted-level
healthy path under each hypothesis. The paths satisfy a rate bound throughout
the complete input history, including every transition slot. The calendar and
terminal horizon are fixed before the noise. Task changes therefore produce
missing-data intervals with continued dynamics. The ideal healthy-subtracted
complete-record information is cubic in the horizon; our question is the
calendar's optimal deficit from that benchmark.

The missing-span loss of the constant-force score classifies this deficit.
A positive loss gives an exact $H^{5/2}$ coefficient and a terminal-balanced
calendar with linearly increasing holds. A zero loss gives $\Theta(H^2)$,
without a sharp quadratic coefficient; its upper construction balances the
complete cumulative task coefficient. Score preservation can occur with a
fixed lower-rank noise image even when affine slope information is lost.
Positive-definite process noise excludes it. The proof covers all legal calendars
in the converse and repairs one scalar mass constraint in the genuine observable
innovation range for attainment. A separate uniform theorem for compact
controller families with common positive-definite noise permits optimization
over both controller and calendar. A fixed bounded task-even bias preserves
the positive coefficient and the quadratic order.

Controlled sensing and controlled Markovian observations already optimize
observation actions for testing \cite{controlled,markov}. Their memoryless or
stationary simple-hypothesis and sequential contracts differ from the present
complete rate-limited paths, growing force, deterministic horizon and missing
state intervals. Monetary switching-cost results also address distinct risks:
Vaidhiyan and Sundaresan take the error limit before sending a sluggish-action
parameter to zero \cite{switchcost}; Lambez and Cohen use known iid laws and
sequential Bayes risk with switching penalty $s=O(c)$ as delay cost $c$ tends to zero
\cite{multipleplays}. Full-history uncertainty sets and constrained active diagnosis
also predate this work: Scott et al. separate complete output sets under bounded
uncertainty \cite{zonotope}, while Gaussian closest-convex-set detectors provide
the fixed-calendar statistical interface used here \cite{convex}.
The present contribution is the second-order calendar deficit within the stated
experiment, rather than a new nearest-set detector or a first-order sensing law.

Feedback design that propagates both mean and covariance is established
\cite{kim}; closed-loop set-valued observers preserve full-history diagnosis
guarantees \cite{closedset}. The feedback/active-diagnosis studies in
\cite{feedback2012,robustfeedback2012} are directly relevant precedents.
Symmetric drift compensation and Allan-variance observing designs account for
elapsed time and real dead time \cite{allan,ossenkopf}, and minimax switchback designs
optimize patterns under fixed horizons \cite{switchback}. Their objectives and
uncertainty contracts differ from ours, but these comparisons do not establish
exhaustive historical priority. Full-text theorem comparisons remain unresolved
for the two 2012 feedback papers and for Cheng--Steinberg's 1991
\emph{Trend robust two-level factorial designs} and Coster--Cheng's 1988
\emph{Minimum Cost Trend-Free Run Orders of Fractional Factorial Designs}.
Their verified bibliographic identities do not establish noncoverage of the
results proved here.

The information restrictions are substantive. The controller has the states
needed for feedback, while the diagnostician reads only the declared retained
states. Fast internal states, commands that reveal them, additional sensors,
or diagnostic readings during transitions define different experiments.
The blackout is a telemetry or admissibility restriction; feedback and
healthy drift continue during it.

\section{The retained-state experiment}\label{sf:model}
Consider the known physical plant and state feedback
\begin{equation}\label{sf:plant}
\begin{aligned}
&x_j=Ax_{j-1}+B_u u_j+B\bigl(\mathbf1_{\{h=1\}}a_j+g_jb_j^h\bigr)\\
&\quad+G\xi_j,\\
&u_j=-Kx_{j-1},\qquad A_K=A-B_uK,\\
&\xi_j\overset{\mathrm{iid}}{\sim}N(0,I_p).
\end{aligned}
\end{equation}
The dimension is fixed, $A_K$ is strictly Schur, $G\in\mathbb R^{d\times p}$
has full column rank, and $0\ne B\in\operatorname{range}G$.
Normality and $\|A_K\|<1$ are unnecessary. Write
$f=G^\dagger B$ and $F=\|f\|^2>0$. The known initial state is
$x_{-n_0}=0$, with fixed $n_0\ge1$. The prefix states at
$j=-n_0+1,\ldots,0$ are all retained, with $g_j=1$ and $a_j=0$.
For $j=1,\ldots,H$, $a_j=\rho_0+\eta j$, where $\rho_0\ge0$ and $\eta>0$.
Throughout, $a_j^h=\mathbf1_{\{h=1\}}a_j$.

Each hypothesis $h\in\{0,1\}$ selects a deterministic complete scalar healthy
path $b^h$ with unrestricted level and adjacent-input rate at most $r/2$,
where $r>0$. Nature observes the fixed deterministic calendar before choosing
these paths and before the independent innovations are drawn. Their actual
difference class is
\begin{equation}\label{sf:path}
\begin{aligned}
&\mathcal Z_H=\{z\in\mathbb R^{H+n_0}:|z_{j+1}-z_j|\le r\\
&\quad\text{ on every adjacent input pair}\}.
\end{aligned}
\end{equation}
Holds retain the complete state in every slot at $g=\pm1$. Changes exclude
at least $k\ge1$ consecutive state readings, with known $|g_j|\le1$ during
the excluded slots. Known exact-$k$ profiles in both directions can be
repeatedly concatenated with holds at any planned start time; they need not
be symmetric and may depend on that predetermined time. Longer transitions
remain legal in the converse. Final holding and retained reads are allowed,
and no terminal task is imposed. No state, noise or healthy level is reset.
There is no hard actuator, state or tracking budget in this analytical class.

Let consecutive retained or initially known states be at $t_{i-1},t_i$,
with $m=t_i-t_{i-1}$. The true predictor is
\begin{equation}\label{sf:predictor}
\begin{aligned}
&x_{t_i}-A_K^m x_{t_{i-1}}\\
&\quad=\sum_{j=t_{i-1}+1}^{t_i}A_K^{t_i-j}\\
&\qquad\bigl[B(a_j^h+g_jb_j^h)+G\xi_j\bigr].
\end{aligned}
\end{equation}
Define its disjoint input map, covariance, and input projector by
\begin{equation}\label{sf:interval}
\begin{aligned}
&H_m=[A_K^{m-1}G,\ldots,G],\\
&W_m=H_mH_m^\top,\\
&P_m=H_m^\top W_m^\dagger H_m .
\end{aligned}
\end{equation}
A gap with $k$ excluded states has $k+1$ inputs, including its right retained
input. Stack the genuine interval projectors to obtain $P_\pi$; inputs after
the last retained state have zero rows and columns in $P_\pi$.
If $U_m$ is an orthonormal basis for $\operatorname{range}W_m$, then
$T_m=(U_m^\top W_mU_m)^{-1/2}U_m^\top H_m$ supplies
$\operatorname{rank}W_m$ nonredundant white rows. The stack $T$ obeys
$TT^\top=I$ and $T^\top T=P_\pi$.

With $\theta=(fa_j)_j$ and $H_gz=(fg_jz_j)_j$, define the profiled information
and its deficit by
\begin{equation}\label{sf:information}
\begin{aligned}
&G_{K,H}(\pi)=\tfrac12\inf_{z\in\mathcal Z_H}\|P_\pi(\theta-H_gz)\|^2,\\
&O_H=\tfrac F2\sum_{j=1}^Ha_j^2,\\
&D_{K,H}(\pi)=O_H-G_{K,H}(\pi),\\
&D_{K,H}^*=O_H-\sup_\pi G_{K,H}(\pi).
\end{aligned}
\end{equation}
The factor $1/2$ is the KL convention. Appendix~\ref{sf:geometry} proves
common support and this exact projection using the actual predictor,
without resetting the states or dropping cross-history covariance.
$O_H$ is the KL of the complete white record after ideal subtraction of
the healthy input. It differs from the composite full-record information,
which still profiles over $z$. Both the complete-record experiment and its
oracle are neutral to known feedback. The supremum in
\eqref{sf:information} need not be attained at finite $H$.

\section{Two deficit regimes}\label{sf:results}
\begin{theorem}[Fixed-plant deficit]\label{sf:two-regime}
Under the experiment of Section~\ref{sf:model}, put
\begin{equation}\label{sf:penalty}
\begin{aligned}
&M_m=\sum_{q=0}^{m-1}A_K^qB,\\
&\beta_\ell=(\ell+1)F\\
&\quad-M_{\ell+1}^\top W_{\ell+1}^\dagger M_{\ell+1}\ge0.
\end{aligned}
\end{equation}
If $\beta_k>0$, then
\begin{equation}\label{sf:positive}
\begin{aligned}
&D_{K,H}^*=\frac{\sqrt2}{5}\sqrt{F\beta_k r}\,\eta^{3/2}H^{5/2}\\
&\quad+o(H^{5/2}).
\end{aligned}
\end{equation}
For every fixed $\xi>0$, the terminal-balanced calendar described below obeys
\begin{equation}\label{sf:upper}
\begin{aligned}
&D_{K,H}(\pi_H^\xi)\le\frac1{10}\\
&\quad\left(\frac{2\beta_k\eta^2}{\sqrt\xi}+Fr\eta\sqrt\xi\right)H^{5/2}\\
&\quad+O_\xi(H^2),
\end{aligned}
\end{equation}
and attains the coefficient in \eqref{sf:positive} at
$\xi_*=2\beta_k\eta/(Fr)$. If $\beta_k=0$, there exist fixed
$0<c<C<\infty$ and $H_0$ such that
\begin{equation}\label{sf:zero}
cH^2\le D_{K,H}^*\le CH^2,\qquad H\ge H_0.
\end{equation}
The first branch is leading-coefficient sharp; the second is order sharp and
asserts no sharp quadratic constant. All constants are for the fixed plant,
noise factor, controller and complete scalar path class.
\end{theorem}

The positive construction uses symmetric blocks with even $n_\ell$:
$n_\ell/2$ positive reads, $k$ missing slots, $n_\ell$ negative reads,
$k$ missing slots, and $n_\ell/2$ positive reads. Their lengths increase as
$n_\ell=\xi\ell+O(1)$; packing leaves one to four final positive reads.
Appendix~\ref{sf:positive-proof} defines the exact input starts, midpoints
and half-slot gap offsets, proves the calendar-uniform converse, and constructs
an observable direction whose single scalar mass is exactly zero.
The zero branch instead uses a profile-aware calendar with bounded cumulative
task coefficient, proved together with its all-calendar lower bound in
Appendix~\ref{sf:zero-proof}.

\begin{proposition}[Preserved score and finite capacity]\label{sf:capacity}
Let $\mathcal R_m=\operatorname{span}\{G,A_KG,\ldots,A_K^{m-1}G\}$,
$\mathcal R_\infty=\bigcup_m\mathcal R_m$, and
$\nu_R=\min\{m:\mathcal R_m=\mathcal R_\infty\}$.
Then $\beta_k=0$ if and only if some $y\in\mathbb R^d$ satisfies
\begin{equation}\label{sf:score}
G^\top(A_K^\top)^q y=f,\qquad q=0,\ldots,k.
\end{equation}
In this case
\begin{equation}\label{sf:capacity-bound}
\begin{aligned}
&k<\nu_R\le\dim\mathcal R_\infty-p+1\\
&\quad\le d-p+1,\qquad k\le d-p.
\end{aligned}
\end{equation}
This preserves the constant scalar force score; it does not require
$P_{k+1}=I_{(k+1)p}$ or recovery of every latent input.
\end{proposition}

\begin{corollary}[Positive-definite process noise]\label{sf:spd}
If $Q=GG^\top\succ0$, then $F=B^\top Q^{-1}B$ and
$\beta_k>0$ for every $k\ge1$. Thus \eqref{sf:positive}--\eqref{sf:upper}
give the fixed-order positive-definite-noise theorem.
\end{corollary}
The score characterization, capacity bound and strict-positive specialization
are proved in Appendix~\ref{sf:geometry}.

\par\smallskip\noindent\textit{Score preservation with lost slope information.}\enspace 
For $k=1$ under the same task protocol, take the nilpotent shift
$A_Ke_2=e_1$, $A_Ke_3=e_2$, $A_Ke_1=0$, with
$G=[e_2,e_3]$ and $B=e_2+e_3$. Then $f=(1,1)^\top$, $F=2$, and
$H_2=[e_1,e_2,e_2,e_3]$ has rank three in a four-dimensional input space.
The constant ray $(1,1,1,1)^\top$ is in its rowspace, so $\beta_1=0$,
but $P_2\ne I_4$. The centered affine ray
$(-1/2,-1/2,1/2,1/2)^\top$ has residual norm squared $1/2$.
Hence preserving the constant score can coexist with loss of slope and other
input modes; the healthy-drift deficit still has quadratic order.

\subsection{Fixed-calendar risk}\label{sf:risk}
The difference class $\mathcal Z_H$ is a finite-dimensional polyhedron, including
its free initial level. Its linear image in the genuine whitened data coordinates
is a closed polyhedron. The Euclidean closest difference therefore exists.
A minimizing $z^*$ realizes both healthy budgets through
$b^0=z^*/2$, $b^1=-z^*/2$. With $\mathcal B_{r/2}$ the complete healthy class,
the true mean sets are $\mathcal M_0=T H_g\mathcal B_{r/2}$ and
$\mathcal M_1=T \theta+T H_g\mathcal B_{r/2}$.
Choose $\mu_0^*=T H_g(z^*/2)$, $\mu_1^*=T \theta-T H_g(z^*/2)$, and put
$d=\|\mu_1^*-\mu_0^*\|=\sqrt{2G_{K,H}(\pi)}$.
For $d>0$, nearest-pair convexity gives
\[
\begin{aligned}
&(\mu_1^*-\mu_0^*)^\top(\mu_0-\mu_0^*)\le0,\\
&(\mu_1^*-\mu_0^*)^\top(\mu_1-\mu_1^*)\ge0
\end{aligned}
\]
for every admissible mean under the respective hypothesis.
The statistic
$S=(\mu_1^*-\mu_0^*)^\top(\widetilde Y-\mu_0^*)/d$
has unit variance and mean at most zero under the null, at least $d$ under
the alternative. Thresholding at $z_{1-\alpha}$ has worst false alarm
at most $\alpha$ and worst power at least
$\Phi(d-z_{1-\alpha})$.
Here $\widetilde Y$ is the stack of whitened true conditional innovations.
The realizing simple pair and the Neyman--Pearson bound show this power is
optimal among tests with the same uniform false-alarm constraint.
Thus, for $0<\alpha,\beta<1$ and $\alpha+\beta<1$, power at least
$1-\beta$ is possible exactly when
\[
G_{K,H}(\pi)\ge\tfrac12(z_{1-\alpha}+z_{1-\beta})^2.
\]
If $d=0$ the closed sets overlap, so no test can have uniform power exceeding
its uniform false-alarm allowance. This is the classical Gaussian closest-set
interface \cite{convex}, not an additional calendar theorem.
\section{Control consequences}\label{sf:control}
When comparing gains, write $\beta_K(\ell)$ for the penalty in
\eqref{sf:penalty} computed with that gain. The original force and noise
are common throughout each controller comparison.
\subsection{Equal poles do not fix diagnostic loss}
Take the same physical plant and process noise
\[
\begin{aligned}
&A=\operatorname{diag}(4/5,3/5),\quad B_u=I_2,\\
&B=(1,2)^\top,\quad Q=\begin{pmatrix}1&1/3\\1/3&5/9\end{pmatrix}.
\end{aligned}
\]
Here $F=29/4$. Gains
$K_d=\operatorname{diag}(3/10,7/20)$ and
$K_n=\left(\begin{smallmatrix}3/10&-2\\0&7/20\end{smallmatrix}\right)$
give closed matrices $\operatorname{diag}(1/2,1/4)$ and
$\left(\begin{smallmatrix}1/2&2\\0&1/4\end{smallmatrix}\right)$.
They share the poles and full oracle, but at $k=1$,
\[
\begin{aligned}
&\beta_{K_d}(1)=1343/344,\\
&\beta_{K_n}(1)=18007/10456.
\end{aligned}
\]
The ratio of sharp deficit coefficients is approximately $0.664$.
This is an ideal known-gain comparison with common $B,Q$ and no actuator
budget; it does not identify a hardware-optimal controller.

\subsection{Settling time changes the controller comparison}
For a scalar common-input plant with $Q=q>0$, $B=1$ and closed parameter
$0\le c<1$, write
\[
\begin{aligned}
&A_m(c)=\sum_{j=0}^{m-1}c^j,\\
&B_m(c)=\sum_{j=0}^{m-1}c^{2j},\\
&\kappa_c(k)=k+1-A_{k+1}(c)^2/B_{k+1}(c).
\end{aligned}
\]
Thus $F=1/q$, $\beta=\kappa/q$ and the coefficient is
$\sqrt{2}\sqrt{\kappa r}\eta^{3/2}/(5q)$.
At fixed $k\ge1$, $\kappa$ strictly decreases in $c\in[0,1)$.
To impose a common homogeneous qualification, let $|e_0|\le E$,
$e_{j+1}=ce_j$, $\gamma=\epsilon/E\in(0,1)$ and a fixed integer
move duration $d_{\rm move}\ge0$. Require the smallest positive $s$ with
$c^s\le\gamma$, and set $k(c)=d_{\rm move}+s(c)$.
This requirement concerns the homogeneous component; forcing and innovations
continue during the excluded slots. It does not assert full noisy return to a tube.

\begin{proposition}[Plateau candidates]\label{sf:settling}
On a compact $0<c_{\rm lo}\le c\le c_{\rm hi}<1$, the minimum of
$\kappa_c(k(c))$ is achieved in the finite set
\[
\begin{aligned}
&\{c_{\rm hi}\}\ \cup\\
&\{\gamma^{1/s}:c_{\rm lo}\le\gamma^{1/s}\le c_{\rm hi},\\
&\qquad\ 1\le s\le s(c_{\rm hi})\}.
\end{aligned}
\]
If $L=-\log\gamma$ and $\delta=-\log c$, then
\[
\begin{aligned}
&\delta\kappa_c(k(c))\longrightarrow L-2\tanh(L/2)>0\\
&\quad(c\uparrow1).
\end{aligned}
\]
\end{proposition}
The proof is in Appendix~\ref{sf:control-proofs}.

For $\gamma=81/100$ and $d_{\rm move}=0$, the finite controller family
$\{4/5,9/10,19/20\}$ has $k=1,2,5$ and exact gap penalties
$1/41$, $2/91$, $5064645/111045881$.
The middle gain is best in that family. On the full interval $[4/5,19/20]$,
the unique best is instead $c=81/100$, with $\kappa=361/16561$.
The latter ranking uses rational isolation of the irrational plateau endpoints.
It optimizes the declared single-gap or coefficient function. Pointwise
calendar sharpness alone does not exchange the limit with optimization over
an entire continuous controller family. For a finite fixed-gain family,
the maximum of finitely many little-oh remainders is still little-oh.


\begin{figure*}[t]
\centering
\IfFileExists{control_settling_figure_v1.pdf}{%
\ifdefined\includegraphics
\includegraphics[width=\textwidth]{control_settling_figure_v1.pdf}
\else
\fbox{\parbox{.94\textwidth}{Analytical figure: the fixed-$k=1$ penalty
decreases with $c$; under the same homogeneous qualification, the penalty
decreases within each settling plateau and jumps upward to its right.
The exact continuum minimum is $c=81/100$, $\kappa=361/16561$.
The companion panel gives the smallest qualifying integer $k(c)$;
equality belongs to the shorter duration.}}
\fi
}{\fbox{\parbox{.94\textwidth}{Analytical figure: the fixed-$k=1$ penalty
decreases with $c$; under the same homogeneous qualification, the penalty
decreases within each settling plateau and jumps upward to its right.
The exact continuum minimum is $c=81/100$, $\kappa=361/16561$.
The companion panel gives the smallest qualifying integer $k(c)$;
equality belongs to the shorter duration.}}}
\caption{Memory benefit after paying the same homogeneous settling requirement. The common scalar plant has innovation variance $q$ and a known force input; $\kappa$ is the dimensionless constant-input missing-span penalty ($\beta=\kappa/q$). For $c\in[.8,.95]$, the qualification uses $E=1$, $\epsilon=.81$, no extra move slots, and the smallest positive integer $s$ with $c^s\le.81$; $k=s$. (a) The fixed-$k=1$ penalty decreases smoothly, whereas the qualified penalty has downward segments and upward right-side jumps. The black square marks the exact continuum minimum at $c=81/100$, $\kappa=361/16561$. (b) The same qualification determines the excluded-slot count. Filled right endpoints belong to the shorter $k$; hollow left endpoints show unattained right limits. Irrational endpoint positions are rounded only for display. The exact three-point family $\{.8,.9,.95\}$ instead has its best member at $.9$; the full-interval ranking comes from rational candidate certificates, not curve resolution. These are analytical curves with no random observations or confidence intervals. The qualification constrains only the homogeneous component and does not certify full noisy/forced return, nonlinear tracking or actuator feasibility.}
\label{sf:settling-figure}
\end{figure*}

\subsection{Joint leading-order design under common positive-definite noise}
\begin{theorem}[Compact common-noise controller family]\label{sf:uniform}
Write $D_H^{*,K}:=D_{K,H}^*$ and
$D_H^{\rm joint}:=O_H-\sup_{K,\pi}G_{K,H}(\pi)=\inf_KD_H^{*,K}$.
Here $\Pi_K$ is the legal calendar class for gain $K$, and the joint supremum
uses $K\in\mathcal K$, $\pi\in\Pi_K$. Let the fixed plant, $B\ne0$,
$Q\succ0$, healthy rate, fault template,
prefix and diagnostic information be common. Let $\mathcal K$ be a nonempty
compact gain set for which each $A_K$ is strictly Schur, and let
$k(K)\in\{1,\ldots,k_{\max}\}$. For each gain assume repeatedly concatenable
known exact-$k(K)$ transitions in both directions, bounded by $|g|\le1$.
Set $C(K)=\sqrt2\sqrt{F\beta_K(k(K))r}\eta^{3/2}/5$. Then
\[
\sup_{K\in\mathcal K}|H^{-5/2}D_H^{*,K}-C(K)|\longrightarrow0,
\]
and $D_H^{\rm joint}/H^{5/2}$ tends to $\inf_K C(K)$.
\end{theorem}
The uniform proof is in Appendix~\ref{sf:control-proofs}.

No continuity of $k(K)$ or of its known profiles is needed for this limit;
neither ensures a finite-horizon optimizer or a minimum of $C(K)$.
In the declared scalar settling example, the finitely many plateau endpoints
are lower semicontinuous because the shorter $k$ is included at each equality.
Its unique coefficient minimum is $c_*=81/100$.
Theorem~\ref{sf:uniform} therefore makes that fixed controller and its
terminal-balanced calendar jointly optimal at leading order under the same
homogeneous qualification. Any sequence of controllers within
$o(H^{5/2})$ of the optimal deficit converges to $c_*$: on the compact set
outside each neighborhood, lower semicontinuity and uniqueness give a positive
coefficient gap, while the information remainder is uniform.
This adds a joint asymptotic conclusion, not a computed finite-$H$ risk optimum
or complete forced-return guarantee.

\section{Bounded task-even uncertainty}\label{sf:even-section}
\begin{corollary}\label{sf:even}
Add $Be_j^h$ with $|e_j^h|\le E/2$ to the input, where $E$ is fixed.
Assume the complete difference class is the explicit Minkowski envelope
$g\mathcal Z_H+\mathcal E_H$, with $0\in\mathcal E_H$ and
$|e_j|\le E$, and leave the covariance, plant and calendar class unchanged.
With $N=H+n_0$ and $A_H=(\sum a_j^2)^{1/2}$, every calendar satisfies
\[
0\le G_{\rm odd}(\pi)-G_{\rm mix}(\pi)
\le FE\sqrt N A_H=O(H^2).
\]
The same bound holds for $D_{\rm mix}^*-D_{\rm odd}^*$ under the same oracle.
Thus the positive leading coefficient is unchanged. Under the fixed
common-support quadratic branch of Theorem~\ref{sf:two-regime}, the same
actual envelope and unchanged projection give
$D_{\rm odd}^*\le D_{\rm mix}^*\le D_{\rm odd}^*+O(H^2)$,
so the order remains quadratic without preserving a quadratic coefficient.
\end{corollary}
The contraction proof is in Appendix~\ref{sf:even-proof}.

This is an order-two perturbation bound, not an exact constant translation.
Multiplicative uncertainty that changes the plant, covariance or support is a
different model. If an unrestricted even difference class contains $a$,
choosing $e=a$ and $z=0$ makes the hypotheses indistinguishable.

\section{Finite verification and a nonlinear risk illustration}\label{sf:verification}
Exact rational checks compare the initialized state covariance against the
retained-input projector in scalar and matrix examples. The matrix report
contains 150 retained-subset checks and 68 full calendars, including rotating,
nonnormal and Jordan models. The proof is universal; these finite checks
support implementation. The symmetric transition profiles in the original
formal reports have zero total repair, so they do not numerically exercise
nonzero $\alpha$. A separately reviewed asymmetric matrix example has
$N=-22863/1307$, $D=628321/10456$ and $\alpha=-182904/628321$.
Its repaired healthy mass is exactly zero. The separately reproduced
six-case asymmetric supplement includes two low-rank examples with repair
energies $100/77$ and $507/178$. These supplement the original zero-repair
checks without changing their historical counts.

The separate planar six-joint mechanical certificate uses a fixed joint-4
fault $f(t)=0.0003(t-1)_+$, four healthy prefix slots, fast step $0.001$ s,
and diagnostic slot $0.25$ s. Move and settle exclude twelve slots while
all process and encoder innovations continue. Point readout is the last
existing pre-step fast sample in the slot; average readout uses its 250 samples.
The readout definitions and rounding were fixed before evaluation.
This comparator has parameter-dependent physical bounds and deterministic
nonlinear approximation errors; it is not claimed to be a direct exact
instance of the fixed common-covariance matrix theorem.


The certificate targets are $\mathrm{PFA}<10^{-3}$ and protected power
at least $0.9$; $H=160$ means 41 seconds including the four prefix slots.
The direction is the frozen nodal projection of the declared terminal-balanced
template, converted to physical task-signed weights and normalized by a
conservative rational norm bound. Exact weights and their KKT/source identities
are supplied with the case certificates; they are chosen before the data.
For every allowed deterministic parameter/path $\vartheta$ under hypothesis $h$,
the fixed physical statistic $Z$ has a coupled unconditional Gaussian comparator
$Z_G^{h,\vartheta}\sim N(\mu_{h,\vartheta},V_{h,\vartheta})$.
The bounds $\mu_{0,\vartheta}\le m_0$, $\mu_{1,\vartheta}\ge m_1$,
$V_{h,\vartheta}\le V^+$, $|Z-Z_G^{h,\vartheta}|\le e_h$ on the good event,
and bad-event probability at most $\delta_h$ hold uniformly over the frozen
instance class. The synthetic fast noises are independent Gaussian modeling
choices: torque variance $10^{-4}$ and encoder variance
$(2\pi/2^{17})^2/12$; they are not measured hardware distributions.
The event can depend on the same innovations: the proof uses event
inclusion and a union bound, not a conditionally Gaussian assumption.
Set
\[
\begin{aligned}
&t=m_0+e_0+\frac{25}{8}\sqrt{V^+},\\
&g_{\rm prot}=m_1-e_1-t.
\end{aligned}
\]
``Protected gap'' denotes $g_{\rm prot}$. With $\overline\Phi$ a conservative upper bound on the standard normal tail,
event inclusion and the unconditional comparator law give
$\mathrm{PFA}\le\overline\Phi(25/8)+\delta_0$: on the null good event,
$Z>t$ implies $Z_G^{0,\vartheta}>t-e_0$.
On the alternative good event, $Z\le t$ implies
$Z_G^{1,\vartheta}\le t+e_1$. When $g_{\rm prot}>0$, protected power is at least
$1-\overline\Phi(g_{\rm prot}/\sqrt{V^+})-\delta_1$, using a conservative
upper normal-tail bound $\overline\Phi$. When the gap is nonpositive this
frozen sufficient certificate returns zero, without substituting the upper
variance into a negative-gap power calculation.
The displayed table values are rounded; exact rational/interval certificate
endpoints, not these display decimals, are used for risk calculation.
For the average case each deterministic allowance is approximately
$0.001148192$ and each dynamic mean loss is $0.000857481$;
for the point case they are $0.023268592$ and $0.009098488$.
The mean endpoints already include the dynamic loss on both sides, and each
bad-event budget is approximately $2.73873\,10^{-11}$.
The numerical certificate records bind the exact direction, both mean/error/event
budgets, and the complete-history covariance witnesses to this fixed case. This presentation does not borrow the
exact common-covariance nearest-set formula to certify the nonlinear plant.

As a representative mechanical case, select the earliest declared
terminal-balanced horizon with at least three complete symmetric blocks and
an average-readout certificate meeting both targets. This gives $H=160$:
three complete blocks, six transfers, and 41 seconds including the prefix.
The block requirement selects this illustration; it is not an extra
requirement for every risk certificate. In particular, the first certified
average-readout grid point remains $H=100$.
For the selected $H=160$ pair, complete-history comparator variance
certificates and two-sided mean/error/event budgets give
\begin{center}
\begin{tabular}{@{}l@{\hspace{4pt}}c@{\hspace{4pt}}c@{\hspace{4pt}}c@{}}
\toprule
Readout & \shortstack{Upper\\variance} & \shortstack{Protected\\gap} & \shortstack{Power lower\\bound}\\
\midrule
\shortstack[l]{250-sample\\average} & $6.1910514\,10^{-7}$ & $0.0290839$ & $0.999999999973$\\
\shortstack[l]{Last existing\\point} & $0.0181313$ & $-0.449714$ & $0$\\
\bottomrule
\end{tabular}
\end{center}
Both false-alarm upper bounds are approximately $0.000967122599$.
The point row fails the sufficient power requirement; its lower bound zero is
not a statement that actual power is zero. The inherited equilibrium enclosure
is restricted to the frozen ramp through $H\le600$
and actual $f\le0.045$, not for every signal in the wider historical force box.
Figure~\ref{sf:readout-figure} extends this finite comparison over the complete
declared terminal-balanced subgrid $H\in\{40,60,80,100,120,160,200\}$.
Every nondegenerate row is retained, including failed sufficient power
certificates. The two $H=40$ directions are degenerate and have no risk bound.
The first certified average-readout grid point for this prescribed family is
$H=100$, or 26 seconds including the prefix. This is neither a sequential
stopping time nor an optimum over calendars.
The declared 400-row family is not completely certified; this closed subgrid
does not complete that calculation.

\begin{figure*}[t]
\centering
\IfFileExists{d2a_readout_figure_v1.pdf}{%
\includegraphics[width=\textwidth]{d2a_readout_figure_v1.pdf}
}{\fbox{\parbox{.94\textwidth}{Finite certificate figure: prescribed
terminal-balanced family, complete declared subgrid through $H=200$.
Average readout first meets the stated grid target at $H=100$;
the point readout's sufficient power bound is zero at all plotted points.
Neither this statement nor the optional illustration asserts actual power
zero, a continuous stopping time or a global optimum.}}}
\caption{Finite certificates for paired readouts of a prescribed
terminal-balanced calendar family. The frozen nonlinear example uses four
healthy prefix slots, 0.25-second slots and the fixed template
$f(t)=0.0003(t-1)_+$. Each horizon has its own declared calendar, shared
by its two readouts. (a) Uniform protected-power lower bounds of the
prescribed statistics; all plotted false-alarm upper bounds are below
$10^{-3}$. The first certified average-readout grid point in this family is
$H=100$. The point bounds are zero because the sufficient certificate gives
no positive guarantee; actual power is not asserted to be zero.
(b) Complete-history Gaussian-comparator variance upper bounds for these
unstandardized statistics. Directions use the prescribed rational norm
scaling, while the statistics are not divided by their standard deviations.
Their weights depend on the horizon and readout; the variance ratio alone
does not quantify universal readout improvement. The true nonlinear
statistics are controlled separately through remainder and event allowances.
The two $H=40$ directions have no risk certificate and are omitted from the
numeric axes. Lines only connect certified grid values; these are bounds,
not measured frequencies or confidence intervals. This finite comparison
does not instantiate the linear asymptotic theorem or complete the
remaining full-grid calculation.}
\label{sf:readout-figure}
\end{figure*}

\section{Limits and conclusions}\label{sf:conclusion}
The two calendar laws require a known fixed stable plant and controller,
complete retained states, a fixed full-column Gaussian noise factor, and a
nonzero scalar fault/healthy force in its image. They do not cover noisy
partial readings, unknown matrices or onset, arbitrary vector nuisance,
saturation or adaptive calendars. Finite healthy amplitude boxes have a
different long-horizon regime. The uniform controller theorem additionally
requires common positive-definite process noise and compactness; near-singular
or rank-changing families require separate analysis. Homogeneous settling is
not a guarantee of full forced or noisy return. The nonlinear illustration
certifies only its stated fixed instance and provides no hardware optimum.

Within this experiment, the constant-force score loss separates a sharp
$H^{5/2}$ deficit from an order-sharp quadratic deficit. The observable
calendar constructions expose how controlled memory, complete healthy drift
and missing time interact. The controller results turn that information law
into a leading-order design rule under common physical inputs and timing.

\appendices
\section{Common support, interval geometry and path duality}\label{sf:geometry}
In the proofs, $P=P_\pi$ denotes the actual input projector,
$\theta=(fa_j)_j$, and $\gamma=(g_jf)_j$. All sums include the fixed prefix
unless a post-input range is specified.
\begin{proof}[Projection, score characterization and capacity]\leavevmode
Let $\mathcal F_{A_K}$ be the initialized causal map from process
inputs to states, and $R$ retain the declared states. Their noise map is
$L_R=R\mathcal F_{A_K}\operatorname{diag}(G)$.
Every mean under either hypothesis is $L_R\vartheta$ with
$\vartheta=(f(a_j^h+g_jb_j^h))_j$. All laws consequently have the
same support $\operatorname{range}L_R$.
For the reduced SVD $L_R=U_s\Sigma_sV_s^\top$, the data on this support
are equivalent to $\Sigma_s^{-1}U_s^\top Y$, with standard Gaussian
noise and mean $V_s^\top\vartheta$. Thus their KL is precisely
$\|V_s^\top(\vartheta_1-\vartheta_0)\|^2/2$.

The predictor transform in \eqref{sf:predictor} is invertible and
block triangular, with identity diagonal blocks. It uses the actual previous
retained state, including the known initial state for the first interval.
It changes $L_R$ into disjoint interval maps $H_m$ and preserves its
rowspace. Therefore $P_m$ is an orthoproject and
\eqref{sf:information} is the exact profiled KL.
This uses rowspace invariance, not an invalid congruence rule for
pseudoinverses. If $U_m$ is an orthonormal basis matrix for $\operatorname{range}W_m$, the
nonredundant white rows are
$(U_m^\top W_mU_m)^{-1/2}U_m^\top H_m$, with identity covariance
of dimension $\operatorname{rank}W_m$.
The full record instead recovers
$G^\dagger(x_j-A_Kx_{j-1})$, giving the stated control-neutral oracle.

We also need a bound valid for all span lengths.
Cayley--Hamilton makes $\mathcal R_d$ invariant and
$\mathcal R_m=\mathcal R_d$ for $m\ge d$.
For such $m$, $W_m\succeq W_d$ with the same kernel, so inverse
monotonicity on $\mathcal R_d$ gives $W_m^\dagger\preceq W_d^\dagger$.
Hence
\begin{equation}\label{sf:gamma}
\begin{aligned}
&\Gamma_K=\max_{1\le m\le d}\|W_m^\dagger\|<\infty,\\
&\sup_{m\ge1}\|W_m^\dagger\|\le\Gamma_K,\\
&S_B=\sum_{q\ge0}\|A_K^qB\|<\infty,\\
&C_K=\Gamma_KS_B^2 .
\end{aligned}
\end{equation}
The last summability follows from strict stability, without normality
or $\|A_K\|<1$. The inverse bound does not use $W_1^\dagger$ on
subspaces reached only later.

With $d_m=\mathbf1_m\otimes f$, $H_md_m=M_m$ and
\[
\beta_{m-1}=\|(I-P_m)d_m\|^2.
\]
Zero loss is exactly $d_m\in\operatorname{row}H_m$, which is
\eqref{sf:score}.
The sequence $\beta_\ell$ is nondecreasing: on a common initialized
horizon of $m+1$ constant-shift inputs, retaining times $1,m+1$
has norm loss $\beta_{m-1}$, while deleting the time-$1$ state has
loss $\beta_m$; rowspace inclusion makes loss nondecreasing.
Also $0\le M_m^\top W_m^\dagger M_m\le C_K$, so
\begin{equation}\label{sf:tail}
\beta_\ell\ge(\ell+1)F-C_K,\qquad \beta_\ell/\ell\longrightarrow F.
\end{equation}

Before reachability stabilizes, its dimension grows by at least one:
$\mathcal R_{m+1}=\mathcal R_m$ already implies invariance.
Since $\dim\mathcal R_1=p$, this proves the upper bounds on $\nu_R$.
If \eqref{sf:score} holds, set $\zeta=(A_K^\top-I)y$.
Then $\zeta\perp\mathcal R_k$.
If $k\ge\nu_R$, this annihilates the invariant reached space, and
$G^\top(A_K^\top)^q y$ has zero successive differences for every
$q\ge0$. It remains the nonzero vector $f$, whereas strict stability
makes it tend to zero. This contradiction proves
\eqref{sf:capacity-bound}.

For the positive-definite specialization, $G$ is square invertible.
If $m\ge2$ and $H_m^\top y=d_m$, its last two blocks give
$G^\top y=f$ and $G^\top A_K^\top y=f$. Thus
$y=Q^{-1}B\ne0$ and $A_K^\top y=y$, contradicting strict stability.
This proves Corollary~\ref{sf:spd} for every nonempty internal gap.
\end{proof}

\par\smallskip\noindent\textit{Affine spans without reflection symmetry.}\enspace 
For a span of length $m$, write $q_j=j-(m+1)/2$,
$d_q=(q_jf)_j$, $N_m=\sum_{j=1}^m q_jA_K^{m-j}B$.
At center amplitude $A\ge0$ its exact oracle norm loss is
\begin{equation}\label{sf:affine}
\begin{aligned}
&\|(I-P_m)(A d_m+\eta d_q)\|^2\\
&\quad=\beta_{m-1}A^2+2\eta A L_{1,m}+\eta^2L_{2,m},\\
&L_{1,m}=-M_m^\top W_m^\dagger N_m,\\
&L_{2,m}=\|(I-P_m)d_q\|^2\ge0 .
\end{aligned}
\end{equation}
By \eqref{sf:gamma}, $|L_{1,m}|\le mC_K/2$.
Thus the negative cross terms over any calendar cost at most
\begin{equation}\label{sf:cross-budget}
E_H=C_K\eta(\rho_0+\eta H)H=O(H^2).
\end{equation}
No reflection symmetry is needed.

The unprojected constant--slope inner product is zero because
$\sum_jq_j=0$. This justifies the stated $L_{1,m}$ even for nonnormal
matrix dynamics; it does not set the projected cross term to zero.

\par\smallskip\noindent\textit{Complete-path support and an observable dual.}\enspace 
For scalar weights on the complete input clock define
$q_t(\lambda)=\sum_{j=-n_0+1}^t\lambda_j$.
Writing each path as its free first level plus adjacent increments and
summing by parts gives
\begin{equation}\label{sf:support}
\begin{aligned}
&h_{\mathcal Z_H}(\lambda)=\\
&\quad\begin{cases}r\displaystyle\sum_{t=-n_0+1}^{H-1}|q_t(\lambda)|,&\displaystyle\sum_j\lambda_j=0,\\+\infty,&\displaystyle\sum_j\lambda_j\ne0.\end{cases}
\end{aligned}
\end{equation}
Indeed the increment at $t+1$ has coefficient $-q_t$, so the independently
allowed increments attain the finite sum by choosing
$\Delta z_{t+1}=-r\operatorname{sign}q_t$. A nonzero total coefficient makes
the free first level unbounded. The same formula on a local block uses
its actual consecutive input slots, including missing ones.

For $Pv=v$ let $\lambda_j=\gamma_j^\top v_j$. On each actual span the
coefficient $c_m=W_m^\dagger H_mv_m$ acting on the predictor innovation has
latent coefficient $H_m^\top c_m=P_mv_m=v_m$, mean $v_m^\top\theta_m$
and variance $\|v_m\|^2$. Hence this range vector is a statistic of retained
states alone, also for singular $W_m$. Applying
$\|u\|^2/2\ge v^\top u-\|v\|^2/2$ before the healthy-path infimum yields
\begin{equation}\label{sf:dual}
\begin{aligned}
&G_{K,H}(\pi)\ge\langle v,\theta\rangle-\tfrac12\|v\|^2\\
&\quad-h_{\mathcal Z_H}(\lambda),\\
&D_{K,H}(\pi)\le\tfrac12\|\theta-v\|^2+h_{\mathcal Z_H}(\lambda)\\
&\quad=\tfrac12\|(I-P)\theta\|^2+\tfrac12\|P\theta-v\|^2\\
&\qquad+h_{\mathcal Z_H}(\lambda).
\end{aligned}
\end{equation}
No data from an unobserved input or fast state enter this dual.
\section{Positive-score-loss converse and attainment}\label{sf:positive-proof}
\begin{proof}[Positive branch of Theorem~\ref{sf:two-regime}]\leavevmode
Assume $\beta_k>0$. Monotonicity and \eqref{sf:tail} imply
$b_K=\inf_{\ell\ge k}\beta_\ell/\ell>0$: with
$L_0=\max\{k,\lceil2C_K/F\rceil\}$, the tail and finite initial part give
$b_K\ge\min\{F/2,\beta_k/L_0\}>0$.
\par\smallskip\noindent\textit{Full-history calendar-uniform converse.}\enspace 
Replacing terminal moves without later readings by holding preserves the old
data and adds readings. Data processing for each full healthy pair and then
an infimum show that the supremum can be restricted to calendars with a
retained state at $H$. Prefix states are all read, so every non-singleton interval is
post-onset and internal.

Fix $\varepsilon\in(0,1)$ and divide $[\lceil\varepsilon H\rceil,H]$
into complete macros
of length $L=\lfloor H^{3/4}\rfloor$.
On each clipped observed run of length $n$, index its nodes by
$i=0,\ldots,n-1$ and use the scalar signed tent
$z_i=g_i r\min(i,n-1-i)$.
Set it to zero on missing inputs, right-first retained inputs, run endpoints,
prefix, macro joins and unused tail. This is one full rate-$r$ path, realized
by $z/2,-z/2$ under the two hypotheses.
Writing $h=gz$, every non-singleton input span has residual $a-h=a$.
If $T$ is the sum of its actual gap norm losses, orthogonal decomposition gives
\begin{equation}\label{sf:converse}
2D_{K,H}(\pi)\ge F\sum_j(2a_jh_j-h_j^2)+T .
\end{equation}
No stationary endpoint or prefix replacement is needed.

In a macro let its minimum amplitude be $A$, its number of missing slots $d$,
its clipped run count $M$ and lengths $n_i$, with $\sum n_i=L-d$.
Put $C=r(2A-rL/2)/4>0$ for large $H$.
For a run, $\sum_{i=0}^{n-1}\min(i,n-1-i)\ge n^2/4-n/2$
and each tent height is at most $rL/2$. Thus its saving is at least
$Cn^2-2Cn$, and the macro saving is at least $C\sum n_i^2-2CL$.
There are $M-1$ complete actual gaps wholly within the macro, each excluding at least
$k$ slots and having center amplitude at least $A$.
Equation \eqref{sf:affine} and monotonicity give
\[
T\ge\sum_{\rm macros}\beta_kA^2(M-1)-E_H.
\]
For each gap its missing amplitudes are at most twice its input-center
amplitude; hence the same $T$ also satisfies
\[
T\ge(b_K/4)\sum_{\rm macros}A^2d-E_H .
\]
Crossing gaps contribute no complete-gap count but each missing node is
assigned once in this second budget.

Fix $\delta\in(0,1)$ and split the same global $T$ before applying its bounds.
Set $W=(1-\delta)\beta_k$.
For $M>0$, Cauchy and AM--GM lower-bound a macro contribution by
\[
\begin{aligned}
&FC(L-d)^2/M+WA^2M+\delta b_KA^2d/4\\
&\quad-WA^2-2FCL\\
&\quad\ge2A\sqrt{FCW}(L-d)+\delta b_KA^2d/4\\
&\qquad-WA^2-2FCL.
\end{aligned}
\]
The condition $A\ge32FWr/(\delta^2b_K^2)$ absorbs its missing-slot term,
uniformly on late macros. For $M=0,d=L$, the dead budget gives the same
bound $2A\sqrt{FCW}L-WA^2-2FCL$ without division by zero.
The cross error is paid once after the convex split.

As $A=\Theta(H)$ and $C=(rA/2)(1+O(H^{-1/4}))$, the leading sum is
$\sqrt{2FrW}\sum A^{3/2}L$.
Coefficient, rounding and incomplete Riemann-block errors are $O(H^{9/4})$,
$\sum A^2=O(H^{9/4})$, $\sum CL=O(H^2)$ and $E_H=O(H^2)$.
All constants are calendar-independent.
The Riemann sum equals
$(2/5)\eta^{3/2}(1-\varepsilon^{5/2})H^{5/2}+o(H^{5/2})$.
Taking the optimal deficit before the limit and then
$\varepsilon,\delta\downarrow0$ proves
\begin{equation}\label{sf:lower}
\liminf D_{K,H}^*/H^{5/2}\ge
\frac{\sqrt2}{5}\sqrt{F\beta_k r}\,\eta^{3/2}.
\end{equation}

\par\smallskip\noindent\textit{Packing and genuine singleton joins.}\enspace 
A symmetric block with even $n=2h$ has $h$ reads at $+$, $k$ missing slots,
$n$ reads at $-$, $k$ missing slots and $h$ reads at $+$.
For block $\ell$ define its length, preceding post-input count,
input midpoint and center amplitude by
\begin{equation}\label{sf:block}
\begin{aligned}
&L_\ell=2(n_\ell+k),\quad s_\ell=\sum_{i<\ell}L_i,\\
&c_\ell=(L_\ell+1)/2=n_\ell+k+\tfrac12,\\
&A_\ell=\rho_0+\eta(s_\ell+c_\ell).
\end{aligned}
\end{equation}
Its inputs are $s_\ell+1,\ldots,s_\ell+L_\ell$. For local scalar weights
write $q_t(\lambda)=\sum_{j=1}^t\lambda_j$, $1\le t<L_\ell$.
With budget $H-1$, fit the maximal number of blocks with
$\bar n_\ell=2\lceil\xi\ell/2\rceil$; evenly distribute remaining multiples of
four slots, adding two to $n_\ell$ per unit and using a fixed left-to-right
tie rule. Observe the last one to four states at $+$.
Every block padding is $O(1)$ and
\[
\begin{aligned}
&M=\sqrt{H/\xi}+O(1),\quad n_\ell=\xi\ell+O(1),\\
&A_\ell=\rho_0+\eta(s_\ell+c_\ell)\\
&\quad=\eta\xi\ell^2+O(\ell+1).
\end{aligned}
\]
At an artificial join both adjacent states are actually read, so the next
conditional input interval is a singleton.
The global input projector therefore splits into its genuine block spans.
It does not reset any state or make output covariances independent.

\par\smallskip\noindent\textit{Vector auxiliary and one scalar weighted constraint.}\enspace 
On active blocks, $A\ge r(k+n-1)/2$, use $h=n/2$, $c_0=(k+1)/2$ and local levels (the positive indices run from $1$ to $h$,
and the negative indices from $1$ to $n$)
\[
\begin{aligned}
&z_i^{+,L}=r(c_0+h-i),\\
&z_i^-=-r[c_0+\min(i-1,n-i)],\\
&z_i^{+,R}=r(c_0+i-1).
\end{aligned}
\]
Held white auxiliary inputs are $U_j=f(A-g_jz_j^{loc})$ and missing ones $fA$.
On inactive blocks, prefix and tail set $U=0$.
Only uniformly finitely many early blocks are inactive.
Local levels generate a dual rather than a claimed global primal path.

The baseline scalar direction has $\lambda_j^0=(g_jf)^\top U_j$ on holds
and zero on missing inputs. It has zero mass and support
\begin{equation}\label{sf:baseline}
h_\ell^0=F\left\{\frac{Ar}{2}n(n+2k)
-4r^2\sum_{j=0}^{n/2-1}(c_0+j)^2\right\}.
\end{equation}
For $q_t(\lambda^0)$, the partial sums have the correct left-positive/right-negative signs,
are zero at the negative center and at the end, and the local levels
saturate $\Delta z=-r\Delta t\,\operatorname{sign}(q_t(\lambda^0))$ wherever $q_t(\lambda^0)\ne0$.
Missing zero weights preserve the physical gap span $k+1$.
The held absolute levels have four copies of
$r(c_0+j)$, $j=0,\ldots,n/2-1$, so
$\sum_{\rm held}g_jz_j^{\rm loc}=rn(n+2k)/2$ and
$\sum_{\rm held}(z_j^{\rm loc})^2=4r^2\sum_{j=0}^{n/2-1}(c_0+j)^2$.
Their saturated support is $\sum\lambda_j^0z_j^{\rm loc}$ by summation
by parts, giving \eqref{sf:baseline}. Here $\Delta t$ is the full input
separation between displayed held nodes; missing-slot interpolation keeps
the same saturated slope along the partial-mass plateau.

Write $\gamma_\ell=(g_jf)_j$, $p_\ell=P_\ell \gamma_\ell$, and
\[
\begin{aligned}
&N_\ell=\gamma_\ell^\top P_\ell U_\ell,\quad D_\ell=\|p_\ell\|^2,\\
&\alpha_\ell=N_\ell/D_\ell,\quad v_\ell=P_\ell U_\ell-\alpha_\ell p_\ell .
\end{aligned}
\]
Apart from two gap-right inputs, the other $2n-2$ held inputs are singletons
with projection $I_p$ and $\|\gamma_j\|^2=F$.
Thus $D_\ell\ge F(2n-2)\ge Fn$.
The difference $(\gamma_j^\top(P_\ell U)_j)_j-\lambda^0$ lies in two fixed
spans only. In each span,
$\sum_j|\gamma_j^\top(PU)_j|\le\sqrt{mF}\|PU\|\le mFA$;
the baseline right-held point contributes at most $FA$.
With $C_k=2(k+2)$ this gives
\[
\begin{aligned}
&|N_\ell|\le C_kFA,\quad |\alpha_\ell|\le C_kA/n,\\
&\gamma_\ell^\top v_\ell=0,\quad P_\ell v_\ell=v_\ell,\\
&\|v_\ell-P_\ell U_\ell\|^2=N_\ell^2/D_\ell\le C_k^2FA^2/n .
\end{aligned}
\]
The repaired mass is one scalar constraint, exactly matching one unrestricted
scalar healthy level. It is not a coordinatewise vector-nuisance constraint.

For $\delta\lambda_j=\gamma_j^\top v_j-\lambda_j^0$,
$\sum\delta\lambda_j=0$ and
$\sum|\gamma_j^\top p_j|\le\|\gamma_\ell\|\|p_\ell\|\le L_\ell F$.
Hence $\|\delta\lambda_\ell\|_1\le C_k(k+3)FA$.
The full scalar rate class has support $r\sum_t|q_t(\lambda)|$ for a zero-mass vector;
each partial sum is at most half its $L_1$ norm.
A length-$L_\ell$ block correction therefore costs
at most $r(L_\ell-1)\|\delta\lambda_\ell\|_1/2=O(FrA(n+k))$.
Summation and restriction of the full path class give
\begin{equation}\label{sf:positive-support}
\begin{aligned}
&h_{\mathcal Z_H}(\lambda)\le\sum h_\ell^0+O(H^2)\\
&\quad\le\frac{Fr}{2}\sum A_\ell n_\ell^2+O(H^2).
\end{aligned}
\end{equation}
The assembled $v$ lies in $\operatorname{range}P$ and
$\sum_j(g_jf)^\top v_j=0$ exactly.

\par\smallskip\noindent\textit{Observable dual and energy costs.}\enspace 
The range condition makes $v$ observable by the predictor pullback in
Appendix~\ref{sf:geometry}. Equation~\eqref{sf:dual} bounds its deficit.
On active inputs, auxiliary scalar errors are at most $(\eta+r)(n_\ell+k)$.
The at most four tail errors are $O(H)$, prefix errors zero, and inactive
blocks contribute $O(1)$.
Therefore $\|\theta-U\|^2=O(H^2)$ and, by contraction and repair,
\[
\begin{aligned}
&\|P\theta-v\|^2\le2\|\theta-U\|^2\\
&\quad+2C_k^2F\sum A_\ell^2/n_\ell=O(H^2).
\end{aligned}
\]
The two input-gap centers lie at offsets $1/2\pm(n_\ell+k)/2$.
With fixed $\beta=\beta_k$ and fixed $L_1,L_2$ from \eqref{sf:affine},
their exact oracle norm loss is
\[
\begin{aligned}
&2\beta(A_\ell+\eta/2)^2+\frac{\beta\eta^2}{2}(n_\ell+k)^2\\
&\quad+4\eta(A_\ell+\eta/2)L_1+2\eta^2L_2 .
\end{aligned}
\]
Consequently $\|(I-P)\theta\|^2=2\beta\sum A_\ell^2+O(H^{3/2})$.
The non-reflection cross has been retained, not asserted zero.

Substituting the support and energy bounds into \eqref{sf:dual} gives
$2D\le2\beta\sum A_\ell^2+Fr\sum A_\ell n_\ell^2+O(H^2)$.
Power summation yields
$\sum A_\ell^2=\eta^2H^{5/2}/(5\sqrt\xi)+O(H^2)$ and
$\sum A_\ell n_\ell^2=\eta\sqrt\xi H^{5/2}/5+O(H^2)$.
This proves \eqref{sf:upper}. Minimizing at
$\xi_*=2\beta\eta/(Fr)$ matches \eqref{sf:lower}, proving \eqref{sf:positive}.
\end{proof}

\section{Score-preserving quadratic deficit}\label{sf:zero-proof}
\begin{proof}[Zero branch of Theorem~\ref{sf:two-regime}]\leavevmode
\par\smallskip\noindent\textit{Terminal dominance.}\enspace  An unread terminal move can be replaced
by continued holding as in Appendix~\ref{sf:positive-proof}; it suffices to
consider calendars that read $H$. Their spans cover every post input.

\par\smallskip\noindent\textit{Zero penalty: a profile-aware upper bound.}\enspace 
Assume $\beta_k=0$ and fix $L>n_0+k+2$.
Keep the full-input cumulative task coefficient
$M_t=\sum_{j=-n_0+1}^t g_j$.
Hold the current sign $s$ until $sM_t\ge L$.
If at least $k+2$ slots remain, execute the known exact-$k$ move,
flip $s$, and continue; otherwise hold to $H$.
This is a deterministic calendar using the known profile at each
planned start. It balances all input coefficients, including the
missing ones, rather than just held signs.

The threshold overshoot is less than one and a move changes $M_t$
by at most $k$. Thus the next opposite hold run reaches its threshold
after between $2L-k$ and $2L+k+2$ slots.
The initial run exceeds two slots; a skipped final move adds at most
$k+1$ holds. For large $H$, every post hold run has at least two
slots and
\begin{equation}\label{sf:balance}
|M_t|\le L+k+2 .
\end{equation}
If $N_h,N_m$ count post held inputs and moves, then
$H=N_h+kN_m$ and $N_h\ge2N_m$.
The singleton count, excluding the right retained point of each gap,
is $S=N_h-N_m\ge H/(k+2)$.

Put $m=k+1$. Since $P_md_m=d_m$, only the affine slope is lost
in each fixed-length gap:
\begin{equation}\label{sf:slope-loss}
\begin{aligned}
&E=\|(I-P)\theta\|^2\le\\
&\quad N_m\eta^2F\,m(m^2-1)/12=O(H).
\end{aligned}
\end{equation}
Let $\gamma=(g_jf)_j$, $w=P\theta$, $p=P\gamma$ and
$\lambda_j=\gamma_j^\top w_j$.
We have $\lambda_j=Fg_ja_j+e_j$, where $e_j=0$ on singletons
and is uniformly bounded on the fixed-length gaps independently
of the amplitude. The post cumulative $S_t=\sum_{j=1}^tg_j=M_t-n_0$
has $|S_t|\le L+k+2+n_0$ by \eqref{sf:balance}. Abel's identity gives
\[
\sum_{j=1}^t g_ja_j=a_tS_t-\eta\sum_{j=1}^{t-1}S_j .
\]
All partial $\lambda$ masses are therefore $O(t+1)$ and
$N=\sum_j\lambda_j=O(H)$. The fixed prefix changes only the constants.

Now $D=\|p\|^2$ lies between $FS\ge FH/(k+2)$ and $F(H+n_0)$.
Set $\alpha=N/D=O(1)$ and $v=w-\alpha p$.
Then $Pv=v$, $\sum_j\gamma_j^\top v_j=0$, and
\begin{equation}\label{sf:global-repair}
\|\theta-v\|^2=E+\alpha^2D=O(H).
\end{equation}
Each $\gamma_j^\top p_j$ is bounded by $F\sqrt m$ because
all nonsingleton spans have fixed length. The repaired scalar weights
$\lambda'_j=\gamma_j^\top v_j$ hence still have partial masses $O(t+1)$.
Equation~\eqref{sf:support} gives the complete-path support bound
\begin{equation}\label{sf:zero-support}
h_{\mathcal Z_H}(\lambda')=
r\sum_{t=-n_0+1}^{H-1}
 \left|\sum_{j=-n_0+1}^t\lambda'_j\right|=O(H^2).
\end{equation}
The free level is handled by the exact repaired zero mass, not by fixing
an initial healthy value.
The observable range vector $v$ satisfies \eqref{sf:dual}, giving
$D_{K,H}(\pi)\le\|\theta-v\|^2/2+h_{\mathcal Z_H}(\lambda')$.
Equation~\eqref{sf:global-repair}, the support bound and \eqref{sf:dual}
prove the $O(H^2)$ upper bound.

\par\smallskip\noindent\textit{Zero penalty: an all-calendar lower bound.}\enspace 
By monotonicity and \eqref{sf:tail}, define finite
\[
\begin{aligned}
&L_z=\max\{m\ge1:\beta_{m-1}=0\}\ge k+1,\\
&b_+=\inf_{m>L_z}\frac{\beta_{m-1}}m>0 .
\end{aligned}
\]
The capacity bound gives $L_z\le\nu_R$.
Set $\epsilon=(4L_z+2)^{-1}$ and $J_0=\lceil\epsilon H\rceil$.
Consider any calendar reading $H$.

If a positive-penalty interval ends at $t\ge J_0$, it lies entirely
in the post affine region and has center amplitude
$A\ge\eta(t+1)/2\ge\eta\epsilon H/2$, even if its span is long.
For sufficiently large $H$, uniformly in the span,
\eqref{sf:affine} gives
$\|(I-P_m)\theta_m\|^2\ge(b_+/2)mA^2$.
Using $z=0$ in the profile therefore gives
\begin{equation}\label{sf:zero-positive-case}
\begin{aligned}
&O_H-G_{K,H}(\pi)\ge\tfrac12\|(I-P)\theta\|^2\\
&\quad\ge \frac{b_+\eta^2\epsilon^2}{16}H^2 .
\end{aligned}
\end{equation}

Otherwise every interval ending at or after $J_0$ has length
at most $L_z$ and preserves $d_m$.
At least $(1-\epsilon)H/L_z-O(1)$ such late intervals cover
inputs $J_0,\ldots,H$. Their center amplitudes are at least
$\eta\epsilon H/2$ for large $H$.
Choose the single complete healthy difference
\[
\begin{aligned}
&z_j=(r/2)g_j,\\
&\widetilde h_j=fg_jz_j=(r/2)fg_j^2 .
\end{aligned}
\]
Because $|g_j-g_{j-1}|\le2$, this path obeys $|\Delta z|\le r$,
and the two paths $z/2,-z/2$ meet their individual rate bounds.
The free prefix level makes the same choice legal there.
For a late zero-penalty span, preservation of $d_m$ gives
\[
\begin{aligned}
&\langle P_m\theta_m,P_m\widetilde h_m\rangle\\
&\quad=\frac r2FA\sum_{j=1}^m g_j^2+\eta\langle P_md_q,P_m\widetilde h_m\rangle\\
&\quad\ge \frac r2FA-C_0 .
\end{aligned}
\]
Here $C_0$ is fixed since $m\le L_z$.
The right retained input has $g^2=1$, so $\sum g_j^2\ge1$.
No entrywise positivity of the projection is assumed.
The total late cross is thus at least
\[
\frac{rF\eta\epsilon(1-\epsilon)}{4L_z}H^2-O(H).
\]
All early intervals end before $J_0$.
Contraction and Cauchy bound their cross from below by
$-(r/2)F\sum_{\rm early}\sqrt{m\sum a_j^2}$.
Their total input count is at most $J_0-1+n_0$, and
$a_j\le\rho_0+\eta J_0$, so the early contribution is at least
$-(r/2)F\eta\epsilon^2H^2-O(H)$.
Finally $\|P\widetilde h\|^2\le r^2F(H+n_0)/4=O(H)$.
Evaluating the profiled infimum at this legal full path and expanding
the half square gives
\begin{equation}\label{sf:zero-preserved-case}
\begin{aligned}
&O_H-G_{K,H}(\pi)\\
&\quad\ge \tfrac12\|(I-P)\theta\|^2+\\
&\qquad \langle P\theta,P\widetilde h\rangle-\tfrac12\|P\widetilde h\|^2\\
&\quad\ge \frac{rF\eta}{2}\left[\frac{\epsilon(1-\epsilon)}{2L_z}-\epsilon^2\right]H^2-O(H)\\
&\quad=\frac{rF\eta\epsilon}{8L_z}H^2-O(H).
\end{aligned}
\end{equation}
The last equality uses $\epsilon=(4L_z+2)^{-1}$.
All constants in both cases are independent of the calendar.
A smaller positive constant than the minimum of the coefficients
in \eqref{sf:zero-positive-case} and \eqref{sf:zero-preserved-case}
gives the lower bound in \eqref{sf:zero}.
Together with the profile-aware upper construction, this completes
the two-regime theorem.
\end{proof}

\section{Controller design proofs}\label{sf:control-proofs}
\subsection{Proof of Proposition~\ref{sf:settling}}
\begin{proof}\leavevmode
For fixed $m=k+1$ the explained constant-mode norm is
$R_m=(1+c)(1-c^m)/[(1-c)(1+c^m)]$.
Its logarithmic derivative is positive because
$\sum_{j=0}^{m-1}c^{2j}>mc^{m-1}$ by strict AM--GM.
The exact increment in $k$ is
\[
\kappa_c(k+1)-\kappa_c(k)
=\frac{(B_m-A_mc^m)^2}{B_m(B_m+c^{2m})}>0.
\]
The plateau $s\ge2$ is $(\gamma^{1/(s-1)},\gamma^{1/s}]$;
the first plateau has right endpoint $\gamma$.
Each nonempty intersection with the compact controller interval contains its
rightmost point. Fixed-$k$ monotonicity places its minimum there.
At the equality $c^s=\gamma$ the shorter $s$ applies, so a floating-point
logarithm must not change the tie rule.
Finally $m\delta=L+O(\delta)$ and
$R_m=\coth(\delta/2)\tanh(m\delta/2)$.
The limit follows, and positivity is
$L-2\tanh(L/2)=\int_0^L\tanh^2(t/2)\,dt>0$.
\end{proof}
\subsection{Proof of Theorem~\ref{sf:uniform}}
\begin{proof}\leavevmode
Continuity of the spectral radius on the compact family gives
$a_*:=\max_K\rho(A_K)<1$. Choose $a_*<\tau<1$.
The resolvent has a finite norm maximum $R_*$ on
$\mathcal K\times\{|z|=\tau\}$. Matrix Cauchy integration gives
$\|A_K^j\|\le R_*\tau^{j+1}$, hence a common finite bound
$S_*\ge\sum_j\|A_K^jB\|$.
Since $W_{m,K}\succeq Q$, the affine-cross constant in Appendix~\ref{sf:positive-proof} is
uniformly bounded by $C_*:=\|Q^{-1}\|S_*^2$.

Each $\beta_K(\ell)$ is continuous and strictly positive for fixed
$1\le\ell\le k_{\max}$. Finite $k$ and compactness give common
$0<\beta_{\rm lo}\le\beta_K(k(K))\le\beta_{\rm hi}$.
The tail $\beta_K(\ell)\ge(\ell+1)F-C_*$ and monotonicity give
a common positive $b_*$ below all $\beta_K(\ell)/\ell$, $\ell\ge1$.
For fixed $\varepsilon,\delta$, the converse proof therefore uses one
absorption threshold and one $o(H^{5/2})$ remainder for all gains and calendars.
Its lower coefficient is
$\sqrt{1-\delta}(1-\varepsilon^{5/2})C(K)$.
First choose $\varepsilon,\delta$ using $\sup C(K)<\infty$, then the common
horizon threshold; this gives the uniform lower comparison.

The optimal $\xi_K=2\beta_K(k(K))\eta/(Fr)$ lies in a fixed positive compact
interval. Packing has a common bounded padding, because the unused time is
less than the next block length and is evenly distributed.
Thus $n_l=\xi_Kl+O(1)$ and $A_l=\eta\xi_Kl^2+O(l+1)$ uniformly.
Because the base total duration is $\xi M(M+1)+O(M)$ uniformly and
the unused budget is less than the next block duration, evenly spreading
its four-slot units adds one common bounded even amount to each $n_l$.
A common finite early cutoff makes every later block active: the center amplitude is at least $\eta\xi_{\rm lo}l(l-1)$, whereas
the active threshold is at most $r(k_{\max}+\xi_{\rm hi}l+C_n)/2$
for one common padding constant $C_n$. The quadratic eventually dominates
the linear bound, giving a common finite inactive prefix. Auxiliary error, repair energy and support correction
are all bounded by common multiples of
$\sum(n_l+k)^3$, $\sum A_l^2/n_l$ and $\sum A_l(n_l+k)$, each $O(H^2)$.
The fixed spans have common affine coefficients, and the two power sums in
Appendix~\ref{sf:positive-proof} have common $O(H^2)$ errors. Hence
$D_H^{*,K}\le C(K)H^{5/2}+C_UH^2$ for one $C_U$.
This proves uniform convergence. Taking infima changes any uniformly bounded
error by no more than its uniform bound, yielding the joint statement.
\end{proof}
\section{Bounded task-even perturbation proof}\label{sf:even-proof}
\begin{proof}\leavevmode
In the real input projection let $t=P(fa)$ and
$\mathcal K=\{P(fgz):z\in\mathcal Z_H\}$.
Each extra mean $P(fe)$ has norm at most $R=E\sqrt{FN}$.
For $d=\operatorname{dist}(t,\mathcal K)$, class inclusion and the triangle
inequality give $(d-R)_+\le d_{\rm mix}\le d$.
Since $0\in\mathcal K$, $d\le\sqrt F A_H$.
For $G=d^2/2$, separately considering $d\ge R$ and $d<R$ gives
$0\le G_{\rm odd}-G_{\rm mix}\le dR\le FE\sqrt N A_H$.
Taking suprema over the same calendars proves the deficit comparison.
\end{proof}
\begin{thebibliography}{13}

\bibitem{controlled}
S. Nitinawarat, G. K. Atia, and V. V. Veeravalli,
``Controlled sensing for multihypothesis testing,''
\emph{IEEE Trans. Autom. Control}, vol. 58, no. 10,
pp. 2451--2464, Oct. 2013, doi: 10.1109/TAC.2013.2261188.

\bibitem{markov}
S. Nitinawarat and V. V. Veeravalli,
``Controlled sensing for sequential multihypothesis testing with controlled
Markovian observations and non-uniform control cost,''
arXiv:1310.1844v2, Jun. 2014.
[Online]. Available: \url{https://arxiv.org/abs/1310.1844v2}

\bibitem{switchcost}
N. K. Vaidhiyan and R. Sundaresan,
``Active search with a cost for switching actions,''
arXiv:1505.02358v1, May 2015.
[Online]. Available: \url{https://arxiv.org/abs/1505.02358v1}

\bibitem{multipleplays}
T. Lambez and K. Cohen,
``Anomaly search with multiple plays under delay and switching costs,''
arXiv:2108.03082v1, Aug. 2021.
[Online]. Available: \url{https://arxiv.org/abs/2108.03082v1}

\bibitem{zonotope}
J. K. Scott, R. Findeisen, R. D. Braatz, and D. M. Raimondo,
``Input design for guaranteed fault diagnosis using zonotopes,''
\emph{Automatica}, vol. 50, no. 6, pp. 1580--1589, Jun. 2014,
doi: 10.1016/j.automatica.2014.03.016.

\bibitem{convex}
A. Goldenshluger, A. Juditsky, and A. Nemirovski,
``Hypothesis testing by convex optimization,''
arXiv:1311.6765v7, Feb. 2016.
[Online]. Available: \url{https://arxiv.org/abs/1311.6765v7}

\bibitem{kim}
K.-K. K. Kim, D. M. Raimondo, and R. D. Braatz,
``Optimum input design for fault detection and diagnosis: Model-based prediction
and statistical distance measures,'' in \emph{Proc. European Control Conf. (ECC)},
Zurich, Switzerland, Jul. 2013, pp. 1940--1945,
doi: 10.23919/ECC.2013.6669785.

\bibitem{closedset}
D. M. Raimondo, G. R. Marseglia, R. D. Braatz, and J. K. Scott,
``Closed-loop input design for guaranteed fault diagnosis using set-valued
observers,'' \emph{Automatica}, vol. 74, pp. 107--117, 2016,
doi: 10.1016/j.automatica.2016.07.033.

\bibitem{feedback2012}
A. Esna Ashari, R. Nikoukhah, and S. L. Campbell,
``Effects of feedback on active fault detection,'' \emph{Automatica},
vol. 48, no. 5, pp. 866--872, May 2012,
doi: 10.1016/j.automatica.2012.02.020.

\bibitem{robustfeedback2012}
A. Esna Ashari, R. Nikoukhah, and S. L. Campbell,
``Active robust fault detection in closed-loop systems: Quadratic optimization
approach,'' \emph{IEEE Trans. Autom. Control}, vol. 57, no. 10,
pp. 2532--2544, Oct. 2012, doi: 10.1109/TAC.2012.2188430.

\bibitem{allan}
R. Schieder and C. Kramer,
``Optimization of radio astronomical observations using Allan variance measurements,''
arXiv:astro-ph/0105071v1, May 2001.
[Online]. Available: \url{https://arxiv.org/abs/astro-ph/0105071v1}

\bibitem{ossenkopf}
V. Ossenkopf,
``The stability of spectroscopic instruments: A unified Allan variance computation scheme,''
arXiv:0712.4335v1, Dec. 2007.
[Online]. Available: \url{https://arxiv.org/abs/0712.4335v1}

\bibitem{switchback}
I. Bojinov, D. Simchi-Levi, and J. Zhao,
``Design and analysis of switchback experiments,''
arXiv:2009.00148v4, Sep. 2025.
[Online]. Available: \url{https://arxiv.org/abs/2009.00148v4}

\end{thebibliography}

\end{document}
